\chapter{VBA Nedir ve Temel Kavramlar}


% ==========================================
% SAYFA 1: GİRİŞ VE MİMARİNİN TEMEL BİLEŞENLERİ
% ==========================================

\section{VBA (Visual Basic for Applications) Nedir?}

VBA, Microsoft Office uygulamalarının arka planında çalışan, iş akışlarını otomatikleştirmek ve özel iş mantıkları (Business Logic) yürütmek için kullanılan nesne yönelimli ve olay güdümlü bir programlama dilidir. Excel başta olmak üzere Word, Outlook, Access gibi tüm Office araçlarında yerleşik olarak bulunur.

\begin{infobox}{\faTerminal\ Projedeki Temel Rolü}
Bu projede VBA; hücre bazlı manuel işlemleri ortadan kaldırarak arka planda çalışan servis fonksiyonları üzerinden vaka tarama, tarih matrisi eşleştirme ve veritabanı kayıt süreçlerini otomatikleştirir.
\end{infobox}

\vspace{0.1cm}

\section{Excel VBA Temel Bileşenleri}

Bir Excel projesindeki VBA bileşenleri varsayılan (yerleşik) ve geliştirici tarafından eklenen özel yapılar olarak ikiye ayrılır. Projede kullanılan mimari bileşenler ve varlık durumları aşağıda özetlenmiştir:

\begin{table}[h!]
\centering
\small
\begin{tabular}{p{0.22\linewidth} p{0.25\linewidth} p{0.46\linewidth}}
\toprule
\textbf{Bileşen Türü} & \textbf{Varlık Durumu} & \textbf{Sistemdeki Görevi ve Kullanım Amacı} \\
\midrule
\textbf{ThisWorkbook} & Yerleşik \textit{(Her Projede var)} & Dosya açılış/kapanış olaylarını dinler.\\
\textbf{Sheet Modülleri} & Yerleşik \textit{(Her Projede var)} & Sayfaya özel olayları dinler. \textit{(Bu projede temiz mimari için boş bırakılmıştır)} \\
\textbf{Standart Modül} & Özel \textit{(Eklenir - \texttt{.bas})} & Arka plan servis fonksiyonlarını ve genel hesaplama mantıklarını barındırır. \\
\textbf{Sınıf Modülü} & Özel \textit{(Eklenir - \texttt{.cls})} & Nesne tabanlı programlama (OOP) sağlar. \\
\textbf{UserForm} & Özel \textit{(Eklenir - \texttt{.frm})} & Form arayüzüdür. \textit{(Bu projede UI karmaşıklığını önlemek için kullanılmamıştır)}. \\
\bottomrule
\end{tabular}
\caption{VBA Bileşenlerinin Projedeki ve Excel'deki Yapısal Durumu}
\end{table}

\begin{warningbox}{\faExclamationCircle\ Proje Mimarisi Notu}
Bu projede kullanıcı arayüzü karmaşıklığını önlemek adına \textbf{UserForm kullanılmamış} ve kod bağımlılığını azaltmak için \textbf{Sheet modüllerine kod yazılmamıştır}. Tüm otomasyon mantığı Standart ve Class modülleri üzerinden servis odaklı yürütülmektedir.
\end{warningbox}

\newpage

% ==========================================
% SAYFA 2: PROGRAMLAMA MANTIĞI VE TEMEL DİL YAPILARI
% ==========================================

\section{VBA Programlama Mantığı ve Temel Yapılar}

VBA (Visual Basic for Applications), nesne tabanlı ve olay güdümlü mimarisiyle performanslı otomasyonlar oluşturulmasını sağlar. Profesyonel VBA geliştirmede bellek yönetimi, değişken bildirimleri ve modüler kod tasarımı temel yapı taşlarıdır.

\begin{tcolorbox}[colback=lightgray,colframe=bordergray,leftrule=3.5pt,boxrule=0.4pt,top=5pt,bottom=5pt]
    \textbf{\textcolor{primary}{\faShield*\ Katı Değişken Tanımlama (\texttt{Option Explicit}):}}
    Modülün en başında yer alan \texttt{Option Explicit} bildirimi, değişkenlerin önceden tanımlanmasını (\texttt{Dim}) zorunlu kılar. Bu standart; yazım hatalarından doğan mantıksal hataları (bug) engeller, yanlış veri türü kullanımının önüne geçer ve bellekte gereksiz alan tahsisini önleyerek çalışma zamanı performansını optimize eder.
\end{tcolorbox}

\begin{tcolorbox}[colback=lightgray,colframe=bordergray,leftrule=3.5pt,boxrule=0.4pt,top=5pt,bottom=5pt]
    \textbf{\textcolor{primary}{\faCube\ Nesne Hiyerarşisi (DOM) ve Bellek Yönetimi:}}
    VBA hiyerarşik bir nesne modeli kullanır (\texttt{Application} $\rightarrow$ \texttt{Workbook} $\rightarrow$ \texttt{Worksheet} $\rightarrow$ \texttt{Range}). Kod içerisinde hücre veya sayfa işlemleri yapılırken nesnelerin bir değişkene atanarak (\texttt{Set ws = ThisWorkbook.Sheets(...)}) başvurulması, Excel'in grafik arayüz yükünü azaltır ve çalışma süresini ciddi oranda kısaltır.
\end{tcolorbox}

\begin{tcolorbox}[colback=lightgray,colframe=bordergray,leftrule=3.5pt,boxrule=0.4pt,top=5pt,bottom=5pt]
    \textbf{\textcolor{primary}{\faCodeBranch\ Prosedür (Sub) ve Fonksiyon (Function) Mimarisi:}}
    VBA'de kod blokları işlevlerine göre iki ana yapıya ayrılır:
    \begin{itemize}
        \item \textbf{Sub (Prosedür):} Belirli bir eylem dizisini çalıştırır ancak geriye bir değer döndürmez. Genellikle buton tetiklemeleri, arayüz ilklendirmeleri ve akış kontrolü için kullanılır (Örn: \texttt{InitializeAppUI}).
        \item \textbf{Function (Fonksiyon):} Parametre alıp veri işleyen ve işlem sonucunda geriye bir değer döndüren yapılardır. Hesaplama, arama ve veri doğrulama işlemleri için tercih edilir (Örn: \texttt{GetReportDate}).
    \end{itemize}
\end{tcolorbox}

\begin{tcolorbox}[colback=lightgray,colframe=bordergray,leftrule=3.5pt,boxrule=0.4pt,top=5pt,bottom=5pt]
    \textbf{\textcolor{primary}{\faDatabase\ Veri Tipleri ve Değişken Seçim Mantığı:}}
    VBA'de performans doğrudan seçilen veri tiplerine bağlıdır. Tam sayılar için \texttt{Long} (16-bit Integer yerine 32-bit bellek optimizasyonu), metinler için \texttt{String}, mantıksal değerler için \texttt{Boolean} ve esnek/bilinmeyen veriler için \texttt{Variant} kullanılır. Yanlış veri tipi seçimi bellek şişmesine veya işlem yavaşlığına yol açar.
\end{tcolorbox}

\vspace{0.1cm}

\begin{warningbox}{\faExclamationTriangle\ Hata Yönetimi ve İstisna Kontrolü (Error Handling)}
VBA'de beklenmeyen hataların uygulamayı çökertmesini önlemek için \texttt{On Error} mekanizmaları kullanılır. Riskli işlemlerden önce \texttt{On Error Resume Next} ile istisnalar es geçilip kontrol edilebilir; işlem tamamlandığında ise sistemin varsayılan hata yakalama moduna dönmesi için hemen \texttt{On Error GoTo 0} ifadesi çalıştırılmalıdır.
\end{warningbox}

\vspace{0.1cm}

\begin{warningbox}{\faTachometer*\ Ekran Yenileme ve Performans Optimizasyonu}
Büyük veri kümeleriyle çalışırken ekranın her işlemde güncellenmesini engellemek için kod başında \texttt{Application.ScreenUpdating = False} ve otomatik hesaplamaları durdurmak için \texttt{Application.Calculation = xlCalculationManual} komutları kullanılır. İşlem bitiminde bu ayarların tekrar aktif edilmesi kritik önem taşır.
\end{warningbox}



\newpage

% ==========================================
% SAYFA 3: ÖRNEK KOD BLOĞU VE DETAYLI İNCELEME
% ==========================================

\section{Projeden Örnek Kod Bloğu ve İncelemesi}

Aşağıdaki kod parçacığı, projenin servis katmanından (\texttt{mod\_Services}) alınmış gerçek bir fonksiyondur. Fonksiyon, kendisine verilen herhangi bir Excel sayfasının ilk satırını tarayarak geçerli rapor tarihini tespit eder:

\begin{lstlisting}[caption={Rapor Tarihi Tespit Servis Fonksiyonu (mod\_Services)}]
' Retrieves the report date from the specified worksheet header row.
Public Function GetReportDate(ws As Worksheet) As String
    Dim col As Long, cellVal As Variant
    GetReportDate = STATUS_NOT_FOUND
    For col = 1 To 50
        cellVal = ws.Cells(1, col).Value
        If Not IsEmpty(cellVal) And IsDate(cellVal) Then
            GetReportDate = Format(CDate(cellVal), "dd.mm.yyyy")
            Exit Function
        End If
    Next col
End Function
\end{lstlisting}

\vspace{0.2cm}
\noindent \textbf{Kodun Çalışma Mantığı ve Adım Adım Analizi}\\

Yukarıdaki kod bloğu, modüler bir VBA servis fonksiyonunun tüm standartlarını içermektedir:

\begin{enumerate}
    \item \textbf{Girdi ve Çıktı Tanımı:} \texttt{Public Function GetReportDate(ws As Worksheet) As String} satırı ile fonksiyon dışarıdan bir Excel sayfası nesnesi (\texttt{ws}) alır ve işlem sonucunda metinsel (\texttt{String}) bir tarih değeri döndürür.
    \item \textbf{Değişken Deklarasyonu:} Döngü sayacı için \texttt{col As Long} ve hücre içeriğini tutmak için esnek veri tipi olan \texttt{cellVal As Variant} tanımlanmıştır.
    \item \textbf{Varsayılan Durum (Fallback):} \texttt{GetReportDate = STATUS\_NOT\_FOUND} ifadesiyle, taramada tarih bulunamazsa fonksiyonun çökmeden güvenli bir durum döndürmesi sağlanır.
    \item \textbf{Döngü Yönetimi ve Doğrulama:} \texttt{For col = 1 To 50} komutu ile ilk satırın 1. sütunundan 50. sütununa kadar olan hücreler taranır. \texttt{If Not IsEmpty(cellVal) And IsDate(cellVal)} koşuluyla hücrenin dolu ve geçerli bir tarih olduğu teyit edilir.
    \item \textbf{Biçimlendirme ve Erken Çıkış:} Tarih bulunduğunda \texttt{Format(..., "dd.mm.yyyy")} ile standart biçime çevrilir ve \texttt{Exit Function} komutuyla kalan sütunlar taranmadan işlem anında bitirilir.
\end{enumerate}

\vspace{0.2cm}

\begin{infobox}{\faSitemap\ Genel Çalışma Mimarisi ve Modülerliğin Avantajı}
Kullanıcı arayüzünde (UI katmanı) yer alan dinamik ListBox, Form ve Buton gibi tüm bileşenler verileri doğrudan hücrelerden okumak yerine bu \textbf{Servis Modüllerini (\texttt{mod\_Services})} çalıştırır. 

Bu modüler yaklaşım sayesinde:
\begin{itemize}
    \item İş mantığı (Business Logic) ile Görsel Arayüz (UI) birbirinden tamamen ayrılır.
    \item Sayfa yapısı veya Excel tasarımı değişse bile sadece servis fonksiyonu güncellenir; arayüz kodlarına dokunmak gerekmez.
    \item Aynı hesaplama ve tarama fonksiyonu projenin farklı yerlerinde tekrar tekrar güvenle kullanılabilir (Reusability).
\end{itemize}
\end{infobox}